\documentclass[11pt,a4paper,fleqn]{article}
\usepackage[T1]{fontenc}
\usepackage{lmodern}
\usepackage{amsmath}
\usepackage[margin=2.5cm]{geometry}
\usepackage[hidelinks]{hyperref}

\setlength{\parindent}{0pt}
\setlength{\mathindent}{2em}
\pagestyle{empty}

\begin{document}

% ===================== Principio portada =====================

\begin{titlepage}
  \centering
  \vspace*{1cm}
  {\Huge\bfseries \underline{Cálculo 1}\par}
  \vfill
  {\huge\bfseries Actividad 1\par}
  \vspace{0.5cm}
  {\Large 23-11-2026\par}
  \vfill
  \begin{tabular}{@{}ll}
    \textbf{Nombre:}      & Adrián Rodríguez Gregorio \\[4pt]
    \textbf{Repositorio:} & \url{https://github.com/AdrianRG90/Calculo_1_entregas} \\
  \end{tabular}
  \vspace*{2cm}
\end{titlepage}

% ======================= Fin portada =======================

\textbf{1)}
\begin{align*}
  \frac{a^{3}b^{4}c^{7}}{a^{-2}b^{5}\sqrt{c}}
    &= \frac{a^{3}b^{4}c^{7}}{a^{-2}b^{5}c^{1/2}} \\
    &= a^{3-(-2)}\, b^{4-5}\, c^{7-1/2} \\
    &= a^{5}\, b^{-1}\, c^{13/2} \\
    &= a^{5}\, b^{-1}\, c^{6}\, c^{1/2} \\
    &= \frac{a^{5}c^{6}\sqrt{c}}{b}
\end{align*}

\textbf{2)}
\begin{align*}
  \frac{2^{3}\,3^{2}}{3^{3}\,2}
    &= \frac{2^{3}}{2^{1}} \cdot \frac{3^{2}}{3^{3}} \\
    &= 2^{3-1} \cdot 3^{2-3} \\
    &= 2^{2} \cdot 3^{-1} \\
    &= 4 \cdot \frac{1}{3} \\
    &= \frac{4}{3}
\end{align*}

\textbf{3)}
\begin{align*}
  \left( \dfrac{\left( 2\dfrac{3}{9} : 3 \right)^{-1}}
               {\left( \dfrac{9}{4} \right)^{2} \left( \dfrac{2}{5} \right)^{-1}} \right)
    &= \dfrac{\left( \left( 2 + \dfrac{1}{3} \right) : 3 \right)^{-1}}
             {\dfrac{81}{16} \cdot \dfrac{5}{2}} \\[6pt]
    &= \dfrac{\left( \dfrac{7}{3} : 3 \right)^{-1}}{\dfrac{405}{32}} \\[6pt]
    &= \dfrac{\left( \dfrac{7}{9} \right)^{-1}}{\dfrac{405}{32}} \\[6pt]
    &= \dfrac{\dfrac{9}{7}}{\dfrac{405}{32}} \\[6pt]
    &= \frac{9 \cdot 32}{7 \cdot 405} \\[6pt]
    &= \frac{32}{315}
\end{align*}

\textbf{4)}
\begin{align*}
  \log_{2} y^{3} &= 6 \\
  3 \log_{2} y   &= 6 \\
  \log_{2} y     &= 2 \\
  y              &= 2^{2} \\
  y              &= 4
\end{align*}

\textbf{5)}
\begin{align*}
  \log_{10} \sqrt[4]{8}
    &= \log_{10} 8^{1/4} \\
    &= \frac{1}{4} \log_{10} 8 \\
    &= \frac{1}{4} \log_{10} 2^{3} \\
    &= \frac{3}{4} \log_{10} 2 \\
    &= \frac{3}{4} \cdot 0.3010 \\
    &= 0.22575
\end{align*}

\textbf{6)}
\begin{align*}
  3\ \text{rad}
    &= 3 \cdot \frac{180^{\circ}}{\pi} = \frac{540^{\circ}}{\pi}
     \approx 171.8873^{\circ} \approx 171^{\circ}\,53^{\prime}\,14^{\prime\prime} \\
    &\quad \left( 0.8873 \cdot 60 = 53.24^{\prime}, \;\;
                  0.24 \cdot 60 \approx 14^{\prime\prime} \right) \\[8pt]
  \frac{2\pi}{5}\ \text{rad}
    &= \frac{2\pi}{5} \cdot \frac{180^{\circ}}{\pi} = \frac{360^{\circ}}{5} = 72^{\circ} \\[8pt]
  \frac{3\pi}{20}\ \text{rad}
    &= \frac{3\pi}{20} \cdot \frac{180^{\circ}}{\pi} = \frac{540^{\circ}}{20} = 27^{\circ}
\end{align*}

\textbf{7)}
\begin{align*}
  \sin\alpha
    &= \sqrt{1 - \cos^{2}\alpha}
     = \sqrt{1 - \frac{1}{16}}
     = \sqrt{\frac{15}{16}}
     = \frac{\sqrt{15}}{4} \\[6pt]
  \tan\alpha
    &= \frac{\sin\alpha}{\cos\alpha}
     = \frac{\sqrt{15}/4}{1/4}
     = \sqrt{15} \\[6pt]
  \sec\alpha
    &= \frac{1}{\cos\alpha} = 4 \\[6pt]
  \csc\alpha
    &= \frac{1}{\sin\alpha}
     = \frac{4}{\sqrt{15}}
     = \frac{4\sqrt{15}}{15} \\[6pt]
  \cot\alpha
    &= \frac{1}{\tan\alpha}
     = \frac{1}{\sqrt{15}}
     = \frac{\sqrt{15}}{15}
\end{align*}

\textbf{8)}
\begin{align*}
  \tan 30^{\circ} &= \frac{h}{271} \\[6pt]
  h &= 271 \cdot \tan 30^{\circ} \\[6pt]
  h &= 271 \cdot \frac{\sqrt{3}}{3} \\[6pt]
  h &\approx 156.46\ \text{m}
\end{align*}

\textbf{9)}
\begin{align*}
  \cos\frac{5\pi}{12}
    &= \cos\left( \frac{\pi}{4} + \frac{\pi}{6} \right) \\[6pt]
    &= \cos\frac{\pi}{4}\cos\frac{\pi}{6} - \sin\frac{\pi}{4}\sin\frac{\pi}{6} \\[6pt]
    &= \frac{\sqrt{2}}{2} \cdot \frac{\sqrt{3}}{2} - \frac{\sqrt{2}}{2} \cdot \frac{1}{2} \\[6pt]
    &= \frac{\sqrt{6} - \sqrt{2}}{4} \\[14pt]
  \cos\frac{7\pi}{6}
    &= \cos\left( \pi + \frac{\pi}{6} \right) \\[6pt]
    &= -\cos\frac{\pi}{6} \\[6pt]
    &= -\frac{\sqrt{3}}{2}
\end{align*}

\textbf{10)}
\begin{align*}
  105^{\circ} &= \frac{210^{\circ}}{2}, \qquad
  \cos 210^{\circ} = -\frac{\sqrt{3}}{2} \\[10pt]
  \sin 105^{\circ}
    &= +\sqrt{\frac{1 - \cos 210^{\circ}}{2}}
     = \sqrt{\frac{1 + \frac{\sqrt{3}}{2}}{2}}
     = \frac{\sqrt{2 + \sqrt{3}}}{2} \\[10pt]
  \cos 105^{\circ}
    &= -\sqrt{\frac{1 + \cos 210^{\circ}}{2}}
     = -\sqrt{\frac{1 - \frac{\sqrt{3}}{2}}{2}}
     = -\frac{\sqrt{2 - \sqrt{3}}}{2} \\[10pt]
  \tan 105^{\circ}
    &= \frac{\sin 105^{\circ}}{\cos 105^{\circ}}
     = -\sqrt{\frac{2 + \sqrt{3}}{2 - \sqrt{3}}}
     = -\sqrt{\frac{(2 + \sqrt{3})^{2}}{4 - 3}}
     = -\left( 2 + \sqrt{3} \right)
\end{align*}

% ================ Principio referencias ================

\newpage
\begin{center}
  \textbf{Referencias}
\end{center}

\begin{list}{}{\setlength{\leftmargin}{1.27cm}\setlength{\itemindent}{-1.27cm}\setlength{\itemsep}{1.2em}}
  \item Anthropic. (2026). \textit{Claude} \url{https://claude.ai}\par
        Se ha utilizado como corrector y como apoyo en el desarrollo y maquetación del formato Latex.
  \item The LaTeX Project. (s.f.). \textit{LaTeX} \url{https://www.latex-project.org}\par
        Se ha utilizado para la redacción de las expresiones matemáticas y la maquetación del documento.
\end{list}

% ================== Fin referencias ==================

\end{document}